\documentclass[11pt]{article}
\usepackage[margin=1.15in]{geometry}
\usepackage{amsmath,amssymb}
\usepackage{microtype}
\usepackage{parskip}
\usepackage[colorlinks=true,linkcolor=black,urlcolor=blue]{hyperref}
\usepackage{booktabs}

\title{What Does a Benign-Trained Anomaly Detector Actually Detect?\\
\large A plain-language account of the core method and findings}
\author{Suramya Angdembay \and Haiyan Tian\\
\normalsize University of Southern Mississippi}
\date{}

\begin{document}
\maketitle

\noindent\emph{This document is a self-contained, simplified companion to
our full paper. It presents only the central method and the central
findings, in linear order, with the vocabulary defined as it appears. The
full paper adds an extensive program of statistical controls and
robustness checks that are omitted here.}

\section{The problem}

Organizations record what employees do on their computers: logons, file
transfers, web activity, e-mail volume, and so on. An \emph{insider
threat} is an employee who misuses legitimate access, for example by
stealing data. Because such events are rare, a common approach is
\emph{anomaly detection}: build a statistical model of normal activity,
and flag days that the model finds unusual.

Formally, a detector is a function $s$ that assigns a real-valued
\emph{anomaly score} $s(x)$ to each observation $x$, where in our setting
$x$ is a record of one user's activity on one day. Higher scores are
meant to indicate more suspicious activity. Detectors are usually
evaluated by how well their scores \emph{rank} malicious days above
benign ones. The standard summary of ranking quality is the following
probability, estimated over a labeled test set:
\begin{equation}
\mathrm{AUC} \;=\; \Pr\bigl( s(X_{\mathrm{mal}}) > s(X_{\mathrm{ben}})
\bigr),
\end{equation}
where $X_{\mathrm{mal}}$ is a randomly chosen malicious day and
$X_{\mathrm{ben}}$ a randomly chosen benign day. (This quantity is
commonly called the ROC-AUC, for ``area under the receiver operating
characteristic curve''; the probability formulation above is equivalent
and is the only form we need.) An AUC of $1$ means perfect ranking and
$0.5$ means the ranking is no better than chance.

The question that motivates our work is simple to state. A high AUC says
that the score \emph{separates} the two classes. It does not say
\emph{on the basis of what information} the separation is achieved. Two
detectors with identical AUC could work very differently: one might
respond to genuinely unusual behavior, while the other might respond to
incidental properties of \emph{who} the malicious users happen to be,
such as their role, department, or personality profile. The second kind
of detector would be nearly useless in practice, because in deployment
the malicious users are not known in advance. Our paper develops a
method for determining which kind of detector one actually has, and
applies it to a modern detector built from a large language model.

\section{The data}

We use two versions (called r4.2 and r6.2) of the CERT insider-threat
corpus, a widely used synthetic dataset that simulates a corporation's
activity logs, produced by Carnegie Mellon University. Each version
contains several thousand simulated employees observed over about
eighteen months, with a small set of employees scripted to carry out
malicious scenarios on particular days. Version r4.2 has sixty malicious
users; version r6.2 has four. We later add a third dataset, TWOS,
collected from real people; it is described in
Section~\ref{sec:findings}.

Each user-day is converted into a short structured text with three kinds
of lines:
\begin{itemize}
\item a \texttt{DAY} line holding static organizational fields for that
user (project, role, business unit, department, team, and an
information-technology administrator flag);
\item a \texttt{PSY} line holding the user's five simulated personality
scores. These are the ``Big Five'' personality traits used in
psychology: openness (O), conscientiousness (C), extraversion (E),
agreeableness (A), and neuroticism (N). In CERT these numbers are
generated by the simulator; they are constant for each user;
\item one \texttt{SES} line per work session, holding numeric behavioral
counts for that session: duration, number of logons, file and
removable-drive activity, e-mail volume, web-browsing category counts,
and similar quantities.
\end{itemize}
A user-day text therefore looks like a small table serialized as text.
The first two lines describe \emph{who the user is} and are essentially
constant across a user's days; the remaining lines describe \emph{what
the user did} that day.

\section{The detector}

\subsection{Language models as probability distributions}

A \emph{language model} is a parametric probability distribution over
sequences of symbols called tokens. Writing a sequence as
$x = (x_1, \dots, x_n)$, the model has the product form
\begin{equation}
p_\theta(x) \;=\; \prod_{t=1}^{n} p_\theta\bigl(x_t \mid x_1, \dots,
x_{t-1}\bigr),
\end{equation}
where each conditional distribution is computed by a neural network with
parameters $\theta$. Modern ``large'' language models are such networks
with billions of parameters, trained on enormous general text corpora.
We use one called Qwen3-8B (eight billion parameters). The network
processes a sequence through $L$ successive layers; at layer $\ell$ and
position $t$ it maintains an internal state vector
$h^{(\ell)}_t \in \mathbb{R}^{d}$ with $d = 4096$. These internal
vectors matter later: they are the objects we will analyze.

\subsection{Adapting the model to normal behavior only}

We specialize the general-purpose model to our corpus by
\emph{fine-tuning}: continuing its training on the serialized user-day
texts, using only days from benign users. The training objective is
maximum likelihood, that is, minimizing the negative log-likelihood of
the benign corpus $\mathcal{D}$:
\begin{equation}
\mathcal{L}(\theta) \;=\; -\sum_{x \in \mathcal{D}} \sum_{t}
\log p_\theta\bigl(x_t \mid x_{<t}\bigr).
\end{equation}
No malicious example is ever used in training. This is the classical
one-class setup: model the normal population, then measure deviation
from it.

For efficiency the fine-tuning does not modify all eight billion
parameters. Each weight matrix $W$ of the network is kept fixed and
receives a trainable low-rank correction,
\begin{equation}
W' \;=\; W + \tfrac{\alpha}{r} BA, \qquad B \in \mathbb{R}^{d_{\mathrm{out}}
\times r},\; A \in \mathbb{R}^{r \times d_{\mathrm{in}}},
\end{equation}
with rank $r = 16$, so only the small matrices $A$ and $B$ are learned.
This technique is called LoRA (low-rank adaptation); we use a
memory-saving variant called QLoRA in which the fixed weights are stored
at reduced numerical precision. The pair (fixed base model, learned
low-rank correction) is called the \emph{adapted model}. The important
point is that the entire effect of training is carried by a
low-dimensional, removable perturbation, which makes ``before versus
after'' comparisons clean.

\subsection{The anomaly score}

The score of a day $x$ is its average surprisal under the adapted
model:
\begin{equation}
s(x) \;=\; -\frac{1}{n}\sum_{t=1}^{n} \log
p_{\theta_{\mathrm{ada}}}\bigl(x_t \mid x_{<t}\bigr).
\end{equation}
Days that the model of normal behavior finds improbable receive high
scores. On the standard evaluation protocol for these benchmarks the
score ranks very well: AUC about $0.95$ on r6.2 and $0.96$ on r4.2. By
the usual standards of the field, this would be reported as a strong
detector. The rest of this document asks what that number is actually
made of.

\section{Looking inside the model}

\subsection{What did fine-tuning change?}

Because the base model is kept frozen, we can run the same day $x$
through both the base and adapted models and subtract their internal
states. Define, at a chosen layer $\ell$ and each token position $t$,
\begin{equation}
\delta_t \;=\; h^{(\ell), \mathrm{ada}}_t - h^{(\ell),
\mathrm{base}}_t \;\in\; \mathbb{R}^{d}.
\end{equation}
The vector $\delta_t$ is precisely the contribution of the fine-tuning
to the model's internal computation at that point. Everything the
detector learned about our corpus, beyond what the general-purpose model
already knew, must be expressed through these vectors.

\subsection{A sparse dictionary for the changes}

The vectors $\delta_t$ live in $\mathbb{R}^{4096}$ and are not
individually interpretable. We therefore approximate them by sparse
nonnegative combinations of learned dictionary elements. Concretely, we
seek a dictionary matrix $D \in \mathbb{R}^{d \times m}$ (with $m = 4d$
columns, called \emph{features}) and, for each $\delta_t$, a coefficient
vector $z_t \in \mathbb{R}^{m}_{\geq 0}$ with at most $k = 8$ nonzero
entries, such that
\begin{equation}
\delta_t \;\approx\; D z_t ,
\end{equation}
with $D$ and the code map chosen to minimize the mean squared
reconstruction error over a large sample of token positions. This is a
sparse coding problem; in the machine-learning literature the particular
parametrization we use (a one-layer encoder network followed by a
keep-the-largest-$k$ rule and a linear decoder) is called a \emph{sparse
autoencoder}. The outcome is that each change vector is expressed as a
short combination of reusable directions, and each direction can be
studied on its own.

\subsection{Selecting candidate features}

Some dictionary features activate similarly on malicious and benign
days; those cannot carry the anomaly signal. We rank features by the
difference in mean activation,
\begin{equation}
g_f \;=\; \mathbb{E}\bigl[z_{t,f} \mid \text{malicious day}\bigr] -
\mathbb{E}\bigl[z_{t,f} \mid \text{benign day}\bigr],
\end{equation}
and take the five features with the largest $g_f$ as the candidate
mechanism. The full paper verifies, through intervention experiments,
that these selected features genuinely influence the anomaly score: if
one edits the coefficients $z_{t,f}$ of the selected features and
propagates the edited internal state through the rest of the network,
the score changes, and it changes more than when the same edit is
applied to matched control features. We omit the details of those
experiments here and take away only their conclusion: the selected
features are causally connected to the score, not merely correlated
with it.

\subsection{Reading off what the features respond to}

Finally we ask \emph{where} the selected features activate within the
day text. Every token position belongs to one line class: the
organizational \texttt{DAY} line, the personality \texttt{PSY} line, or
a behavioral \texttt{SES} line. For a feature $f$, define its
\emph{attribution mass} on class $c$ as the fraction of its total
activation that occurs at positions of that class:
\begin{equation}
m_f(c) \;=\;
\frac{\sum_{(x,t)\,:\, \mathrm{class}(t) = c} z_{t,f}(x)}
     {\sum_{(x,t)} z_{t,f}(x)} .
\end{equation}
If a feature has nearly all of its mass on \texttt{PSY} and \texttt{DAY}
positions, it responds to who the user is; if its mass lies on
\texttt{SES} positions, it responds to what the user did. This single
quantity is what turns the analysis into an answer to our opening
question.

\section{Findings}
\label{sec:findings}

\subsection{Two detectors that look alike and are not}

The two CERT versions produce nearly identical headline AUC values, yet
the analysis above reveals qualitatively different detectors.

On r6.2, the five selected features place between $99.8$ and $100$
percent of their attribution mass on the personality and organizational
lines. The detector's causally relevant evidence is almost entirely a
description of \emph{who the user is}. On r4.2, four of the five
selected features concentrate on the behavioral session lines. There,
the evidence is mostly \emph{what the user did}.

\subsection{The consequence for the headline numbers}

Why would reading a user's static profile produce a high AUC at all? The
adapted model was trained on roughly ninety percent of the benign users.
For those users, the model has memorized the constant profile lines, so
their days look unsurprising. The malicious users, however, were never
seen in training, so their profile lines alone make their days look
surprising. The score is then partly measuring \emph{familiarity} rather
than behavior.

This suggests a sharper test: compare malicious days only against benign
days of users who were also excluded from training, so that both sides
are equally unfamiliar. Under this comparison the AUC falls from $0.95$
to $0.55$ on r6.2, which is barely better than chance, and from $0.96$
to $0.67$ on r4.2. The apparently excellent detector on r6.2 was, to a
first approximation, an unfamiliar-user detector. Consistently with
this, deleting the personality lines from the inputs at scoring time
\emph{improves} the r4.2 detector's performance on unseen users from
$0.67$ to about $0.80$: the profile text was not merely unhelpful, it
was actively obscuring the behavioral signal.

\subsection{A real-data test with real personality scores}

The CERT data are simulated, including the personality scores, so we
repeated the entire pipeline on TWOS, a dataset collected from
twenty-four real users over one week, during which some users, by
design of the underlying study, temporarily took over other users'
sessions. Each real user completed a standard fifty-item Big Five
personality questionnaire, so the serialized days contain genuine
psychometric information. In this dataset half of the users experience
an attack, so ``being an unusual person'' cannot separate the classes,
and our account predicts that the detector must rely on behavior. That
is what the analysis finds: the selected features place essentially all
of their mass on the behavioral lines, under two independent training
runs, and the score meaningfully separates each user's attacked windows
from that same user's normal windows (average within-user AUC about
$0.7$, with identity held fixed by construction).

\subsection{Where the shortcut comes from}

One might suspect that the r6.2 data simply force any statistical method
to use profile information. This is testable. We trained three classical
anomaly detectors (an isolation forest, a one-class support vector
machine, and a principal-component reconstruction method; their details
do not matter here) on the same days, represented as ordinary numeric
feature vectors that \emph{include} the profile variables. All three
essentially ignore the profile variables: those variables receive about
four to seven percent of the total importance, no profile variable
appears among any method's top five, and deleting the profile variables
does not hurt their performance.

The shortcut is therefore specific to the language-model pathway, and
the reason is visible in the training objective. For a
likelihood-based sequence model, the per-user-constant profile tokens
are the most predictable part of every training sequence, so memorizing
them yields the largest and cheapest reduction of the training loss. In
a plain feature-vector representation, the same information is twelve
columns among two hundred forty-two, with no comparable incentive
attached. Our summary of the whole phenomenon is a two-factor statement:
the population structure of a dataset determines \emph{when} an
identity shortcut is available, and the training objective together with
the data representation determines \emph{which} detectors take it.

\section{Why this matters}

Three conclusions carry beyond these particular datasets and this
particular model.

First, for anomaly detectors intended to detect behavior, ranking
quality alone is not evidence of success. A detector should be evaluated
with the unfamiliarity confound removed, by comparing malicious
observations only against normal observations from equally unseen
users, and, where possible, by within-user comparisons.

Second, including static descriptors of people (role, department,
personality measures) in a detector's input is not harmless. In the
likelihood-based detector it created the shortcut, and even where the
detector was behavioral it degraded performance.

Third, the question ``what does this detector actually detect?'' can be
answered, not merely asked. The chain presented here, namely
differences of internal states, a sparse dictionary approximation,
selection of score-relevant directions, and attribution of those
directions to input content, turns the question into a measurable
quantity. In the full paper this chain is supported by causal
intervention experiments and an extensive robustness program, and it is
that combination, rather than any single number, that justifies the
conclusions summarized here.

\end{document}
